\chapter{Methods}\label{ch:methods}

\section{National time series}
\subsection{Waves and segmentation}\label{sec:waves}
Waves are labelled as the four largest peaks of the smoothed daily case series that lie at least 90 days apart. The rule fixes where the peaks are. That there are four is my reading of the curve, and the rule does not estimate it.

The segmentation uses the Pruned Exact Linear Time (PELT) algorithm \parencite{killick2012}, which returns the segmentation that minimises a penalised cost. I apply it to the log of weekly mean cases, allowing the mean and the variance to change, with the modified BIC penalty and a minimum segment length of six weeks. Weekly means are used because the smoothed daily series is so strongly autocorrelated that it over-segments. PELT assumes independent observations within a segment, which an epidemic curve does not satisfy, so I treat its change points as a summary of the shape of the curve that carries no significance level. To show how far they depend on the settings, I repeated the segmentation under two other penalties and minimum segment lengths of four, six and eight weeks, and I report the range of the number of change points.

\subsection{Mobility retreat index}
The six Google mobility series move together. I summarise them by the first principal component of the standardised daily series \parencite{hastie2009}, signed so that higher values mean more time at home and fewer trips out, and call it the \emph{retreat index}.

\subsection{Stringency and mobility}\label{sec:policy-model}
Daily values were averaged to calendar weeks. A regression of the weekly retreat index on stringency in levels, with year-specific intercepts and slopes, leaves residuals with a lag-1 autocorrelation of \PolLevRho{} (Durbin--Watson statistic \PolLevDw). Both series are persistent, and with residuals this strongly autocorrelated the levels regression is at risk of being spurious. A HAC correction with a few lags does not repair that. Inference therefore rests on week-on-week changes,
\begin{equation}
\Delta R_w = \alpha + \beta\, \Delta S_w + \sum_{k\in\{2021,2022\}} \big(\gamma_k D_{kw} + \delta_k\, \Delta S_w D_{kw}\big) + \varepsilon_w,
\end{equation}
where $\Delta R_w$ and $\Delta S_w$ are the changes in the retreat index and the stringency index from week $w-1$ to week $w$, and $D_{kw}$ indicates year $k$. The coefficient $\beta$ is the change in the retreat index that accompanied a one-point change in stringency in 2020, and $\delta_k$ measures how that response differed in year $k$. The residual lag-1 autocorrelation of this model is \PolDifRho. Standard errors are Newey--West with four lags \parencite{newey1987}, which allows for the short-run serial correlation and heteroskedasticity that remain after differencing. As a check, the model was refitted without March 2020, when the first lockdown moved both series at once.

I also regressed next-week log growth in cases on the retreat index lagged by zero to four weeks, with the same standard errors. Restrictions tighten and people stay home when cases are already rising, so this regression cannot separate the effect of mobility on transmission from the reverse, and I report it only briefly.

\section{Household panel}
\subsection{Attrition and paired change}
Women retained at follow-up were compared with those lost on baseline covariates, using Welch $t$-tests and the standardised mean difference. An absolute value below 0.1 is conventionally treated as negligible.

The change in the number of food groups between rounds was tested with the paired $t$-test and the Wilcoxon signed-rank test, with Cohen's $d_z$ (mean paired difference over its standard deviation) as the effect size. For the binary outcomes, MDD-W and each of the ten food groups, McNemar's test compares the numbers of women who gained and who lost the outcome \parencite{mcnemar1947}. The ten food-group $p$-values were adjusted by the Benjamini--Hochberg procedure \parencite{benjamini1995}.

\subsection{Differences by wealth}\label{sec:wealth-models}
Whether the change differed by baseline wealth was examined with two models, fitted to the same \PanAncN{} women with complete covariates. Write $Y_{i1}$ and $Y_{i2}$ for the number of food groups woman $i$ consumed in 2019 and 2020, and $\mathbf{q}_i$ for her wealth-quintile indicators:
\begin{align}
Y_{i2} - Y_{i1} &= \alpha_0 + \mathbf{q}_i^\top\bm\theta + \alpha_1\,\text{district}_i + e_i, \label{eq:change}\\
Y_{i2} &= \beta_0 + \beta_1 Y_{i1} + \mathbf{q}_i^\top\bm\gamma_q + \mathbf{x}_i^\top\bm\gamma_x + \varepsilon_i. \label{eq:ancova}
\end{align}
In the change-score model~\eqref{eq:change} the hypothesis that the gain was the same in every quintile, $\bm\theta = \mathbf{0}$, is tested with a joint Wald test. In the ANCOVA~\eqref{eq:ancova}, $\mathbf{x}_i$ holds district, age, education above primary level, the number of children in the household, the two 2020 constraint indicators, whether the husband's income fell, whether the woman started or stopped outside work, and the change in her food-decision index. Both models use HC3 heteroskedasticity-consistent standard errors.

When groups differ at baseline the two models answer different questions and can disagree, which is Lord's paradox \parencite{lord1967}. Model~\eqref{eq:change} asks whether richer women gained more. Model~\eqref{eq:ancova} asks whether, among women with the same 2019 score, richer women had more diverse diets in 2020. A single 24-hour recall is a noisy measure of usual diet, so $\beta_1$ can be expected to fall well below 1, and then the ANCOVA adjusts little for baseline and its wealth coefficients largely reproduce the gap that already existed in 2019. I use the change-score model as the test of unequal gains and the ANCOVA as a description of 2020 diets.

\section{Mental-health survey}
\subsection{Reliability and prevalence}
Internal consistency of each scale is summarised by Cronbach's $\alpha$ \parencite{cronbach1951} and by McDonald's $\omega_t$ from a one-factor model. Prevalences are reported with 95\% Wilson score intervals \parencite{wilson1927}, whose coverage stays close to nominal in the smaller subgroups.

\subsection{Regression models}\label{sec:mh-methods}
Probable depression was modelled by two nested logistic regressions. Model A contains the sociodemographic and COVID-19 exposure variables, and model B adds the optimism, pessimism and mindfulness scores. The two are compared by a likelihood-ratio test and Nagelkerke's $R^2$. Probable anxiety was fitted with model B only. The three scales were measured in the same questionnaire as the outcomes and may lie on the path between an exposure and symptoms, so model B can adjust away part of an exposure's association. Estimates from model A are therefore shown beside those from model B for the main predictors. Within each model the $p$-values were adjusted by the Benjamini--Hochberg procedure and are reported as $q$.

Nearly half the sample meets the depression threshold. At that prevalence an odds ratio is much further from 1 than the corresponding ratio of prevalences, and it is easily misread as one. I therefore also report adjusted prevalence ratios from modified Poisson regression \parencite{zou2004}, a log-link Poisson model for the binary outcome with sandwich standard errors. Model B was refitted without the flagged records of Section~\ref{sec:suspect} as a sensitivity analysis.

\subsection{Latent profile analysis}
The five standardised scale totals for anxiety, depression, stress, optimism and mindfulness were modelled as a finite Gaussian mixture, $f(\mathbf{z}_i) = \sum_{g=1}^{G} \pi_g\, \phi(\mathbf{z}_i; \bm\mu_g, \bm\Sigma_g)$, fitted by the EM algorithm in \texttt{mclust} \parencite{scrucca2016}. Pessimism was left out because of its low reliability. Latent profile analysis assumes that the indicators are independent within a class, so only the four diagonal covariance parameterisations (EEI, VEI, EVI, VVI) were considered, for $G = 1$ to 8. The model was chosen by ICL \parencite{biernacki2000}, which penalises BIC for overlap between classes, and BIC is reported as well. The models were fitted to the \GmmN{} unflagged respondents, and I call the classes of the selected model profiles. To see what the flagged records do to a mixture model, I also fitted every diagonal model with $G = 3$ to 6 to the full sample and counted the fits that contain a class holding under 5\% of respondents, at least 80\% of them flagged.

Relative entropy and the average posterior probability of each class measure how confidently the fitted model assigns respondents. They say little about whether the classes are separated in the data, so I added checks that do. The mean silhouette width \parencite{rousseeuw1987} was computed for the mixture classification and for $k$-means solutions with $k = 2$ to 8. The mixture classification was compared with a $k$-means partition with the same number of clusters by the adjusted Rand index. The mixture was refitted without the depression indicator, because profiles defined partly by the PHQ-9 must differ in depression prevalence. The first principal component of the five indicators was regressed on profile membership, where a high $R^2$ means that the profiles are bands along one severity dimension. Profile membership was then related to covariates by multinomial logistic regression with the lowest-symptom profile as reference, giving relative risk ratios with Benjamini--Hochberg $q$-values.

\section{Prediction of depressive-symptom status}\label{sec:ml-design}
\subsection{Question and feature sets}
The outcome is probable depression (PHQ-9 $\geq$ 10), and three nested feature sets were used. Tier A has the 15 sociodemographic and COVID-19 exposure variables. Tier B (18 variables) adds the optimism, pessimism and mindfulness scores. Tier C (20) adds the anxiety (BAI) and stress (PSS-4) totals. Tier C contains concurrent symptom scores, and anxiety and depression co-occur, so it measures overlap between symptom scales and is included as a reference. Features are named in an explicit list (Appendix~\ref{app:repro}), which excludes the PHQ-9 items and total and the original authors' derived columns. As a ceiling, I also score the PHQ-2, the sum of the first two PHQ-9 items, on its own. It is part of the outcome and is never a model feature.

\subsection{Learners}
Four learners were compared. Logistic regression was fitted by unpenalised maximum likelihood. The elastic net \parencite{zou2005,friedman2010} was tuned over mixing parameters $\alpha \in \{0, 0.25, 0.5, 0.75, 1\}$ and the penalty $\lambda$ by cross-validated binomial deviance. Deviance was used instead of AUC because AUC is unchanged by any monotone rescaling of the predictions, so tuning on it does not penalise over-shrinkage and can damage calibration. The random forest \parencite{breiman2001,wright2017} has 1,000 probability trees, with the number of candidate variables per split (2, 3, 5 or 8) and the minimum node size (5, 15 or 30) chosen by out-of-bag AUC. XGBoost \parencite{chen2016} used a learning rate of 0.03 and row and column subsampling of 0.8, with tree depth (2, 3 or 4) and minimum child weight (1 or 5) chosen by cross-validated AUC and the number of rounds by early stopping.

\subsection{Validation design}
\paragraph{Fixed hold-out.} The data were split once, stratified by outcome, into a training set (75\%, $n = \MlTrainN$) and a test set (25\%, $n = \MlTestN$). All tuning used the training set only, with the same stratified ten folds for logistic regression, the elastic net and XGBoost, and out-of-bag predictions for the random forest. This split supports the detailed reporting, which covers test AUC with DeLong 95\% intervals \parencite{delong1988,robin2011}, the Brier score, the calibration slope \parencite{vancalster2019}, decile calibration plots, and sensitivity and specificity at the threshold that maximises Youden's $J$ \parencite{youden1950} in the training out-of-fold or out-of-bag predictions. The calibration slope is the coefficient from a logistic regression of the outcome on the logit of the predicted probability. Its ideal value is 1, and a value below 1 means the predictions are too extreme. Where one learner per tier is singled out, it is the learner with the highest cross-validated AUC on the training set (out-of-bag AUC for the random forest), and test AUC plays no part in that choice. Paired DeLong tests compare models on the fixed test set.

\paragraph{Repeated splits.} One test set of \MlTestN{} respondents gives a noisy AUC, and a ranking of learners read from it could be an accident of the split. The comparison between learners therefore rests on \RepN{} repeated stratified 75/25 splits. In each split all twelve models are tuned again on the training part and scored on the test part. To keep the computation manageable, the repeated splits use five-fold cross-validation and smaller tuning grids than the fixed split. I report the mean and standard deviation of test AUC across splits, the within-split difference between each learner and logistic regression, and the share of splits in which the learner beats logistic regression.

\paragraph{Feature ablation.} On the same \RepN{} splits, logistic regression was fitted to nine feature sets: tier A, tier A without division, tier A plus each of the three scales in turn, tier B, mindfulness alone, optimism alone, and the PHQ-2 reference. Division is removed in one of them because the flagged records are concentrated in one division, so part of what division predicts is an artefact of those records.

\subsection{Interpretation and checks}
Permutation importance was computed for the logistic regression models of tiers A and B on the fixed test set. It is the mean drop in test AUC over 30 random permutations of one variable, with all dummy columns of a categorical variable permuted together. The logistic models are the ones interpreted because no learner improved on them in the repeated-split comparison.

The tier-B test AUC was also recomputed after removing the flagged records from the test set, and probable anxiety was predicted from tiers A and B as a secondary outcome. As an implementation check, the tier-B design matrices and the fixed split were exported, and L2-penalised logistic regression, a random forest and histogram gradient boosting were fitted in Python with scikit-learn \parencite{pedregosa2011}, each tuned by stratified ten-fold grid search. The Python models see the same data and the same split, so agreement with R checks the implementation and nothing else.
